\section{Conclusion}
\label{sec:conclusion}

We explored the phase diagram of the extended SSH spin model through a combined correlation diagnostic accessible in a single $Z$ measurement basis.
The staggered structure factor identifies AFM correlations, while the cell-string correlation supplies complementary information on the trivial and topological dimerized regimes.
Their reference-state-calibrated combination $Q$ captures a three-region structure similar to the partial-reflection marker in the $L=8$ numerical map.
A two-site-cell symmetry description connects the dimer limits to distinct projective representations, and the sector-gap trend along $\delta=-2$ is consistent with the gapless uniform-chain prediction.

We prepared states with phase-matched initial states and a reflection-symmetric variational circuit, optimized classically and executed on the Baihua processor.
The experimental protocol combines cluster-state-based qubit screening, tensor-product readout correction, and measurements of both constituent correlators from the same computational-basis samples.
For the displayed data, the $30$-point experiment--simulation comparison uses $16\,384$ shots per point and gives $r=0.86$ with fitted marker scales.

The measurement strategy relies on site-resolved spin or equivalent occupation readout.
Optical-lattice imaging has already accessed nonlocal string correlations~\cite{endres2011observation}, and optical-tweezer Rydberg arrays provide site-resolved readout in a bosonic SSH realization~\cite{deleseleuc2019observation}.
These capabilities motivate transferring the correlation diagnostic to such simulators when the corresponding spin encoding and reference-state calibration are available.
